% Propietario conservado: /Users/ruben/Documents/ChatGPT/jueces y controles/REVISION_ARGUMENTAL_APERTURAS_CIERRES_20260915/II/manuscript_es/sections/07b_geometria_elipse.tex
% Adición separada: no sustituye 07_elipse.tex.
% Procedencia: integral 2249, 52.4.1--2, 52.6.2--3, 53.28 y 90.4.6.
% Recuperación de resultados y explicitación geométrica; certificado local
% verificar_elipse_II.py. No usa un radio físico como entrada generadora.

\subsection{Demi-axes, foyers et rayons intérieurs}
\label{ii:sec:ii-elipse-geometria}

La figure est obtenue à partir des sorties angulaires et d’action des sections précédentes. APP conserve les feuillets ainsi que leurs résidus et quotients ; TRIT fixe l’orientation ; TPK transporte la frontière, la retenue et la mémoire jusqu’aux expressions du même état enrichi et de la structure discrète conjointe du continu. On reçoit ici le couple \((A,C^*)\) et la section positive \(\hbar\) après cette construction. La géométrie explicite l’information conservée par leurs lecteurs, sans utiliser une longueur observée pour sélectionner l’état source. La section~\ref{ii:sec:ii-elipse-geometria} développe cette réalisation à partir des lecteurs déjà engendrés, et la construction angulaire est conservée dans la section~\ref{v:v:ant:sec:ii-angulos}.

Soit \(\xi=C^*/A\in(-1,1)\), avec \(A\ne0\) et les deux angles dans la
même unité. Écrivons \(\eta=\operatorname{artanh}\xi\) et
\[
 a_S=\sqrt{2\hbar}\,e^{\eta/2},\qquad
 b_S=\sqrt{2\hbar}\,e^{-\eta/2},\qquad
 \gamma(\theta)=(a_S\cos\theta,b_S\sin\theta).
\]
Les demi-axes sont étiquetés par leurs coordonnées : si \(\eta<0\),
l'axe vertical est le plus grand. Tous deux appartiennent à la droite des racines carrées
d'action de la carte équilibrée, non à la droite des longueurs spatiales.

\begin{proposition}[Géométrie focale et intérieur radial]
\label{ii:prop:ii-elipse-focos}
Soient \(a=\max(a_S,b_S)\), \(b=\min(a_S,b_S)\). La demi-distance
focale et l'excentricité sont
\begin{equation}
 f_S=\sqrt{a^2-b^2}
 =\sqrt{4\hbar\sinh|\eta|},\qquad
 e_f=\frac{f_S}{a}=\sqrt{1-e^{-2|\eta|}}.
 \label{ii:eq:ii-elipse-focal}
\end{equation}
Pour \(\eta>0\), les foyers sont \((\pm f_S,0)\) ; pour
\(\eta<0\), ils sont \((0,\pm f_S)\). En \(\eta=0\), ils coïncident avec
le centre. Le rayon de frontière dans la direction polaire \(\phi\) est
\begin{equation}
 \rho_{\partial}(\phi)
 =\left(\frac{\cos^2\phi}{a_S^2}
        +\frac{\sin^2\phi}{b_S^2}\right)^{-1/2}.
 \label{ii:eq:ii-elipse-radio-polar}
\end{equation}
Le segment intérieur dans cette direction est constitué des points
\(r(\cos\phi,\sin\phi)\), avec \(0\le r<\rho_{\partial}(\phi)\).
En particulier, \(b\le\rho_{\partial}(\phi)\le a\).
\end{proposition}

\begin{proof}
Les carrés du grand et du petit demi-axe sont
\(2\hbar e^{|\eta|}\) et \(2\hbar e^{-|\eta|}\). Leur différence
donne \eqref{ii:eq:ii-elipse-focal}. Pour vérifier le foyer, supposons
\(\eta>0\). En \(X=\gamma(\theta)\), les distances à
\(F_+=(f_S,0)\) et \(F_-=(-f_S,0)\) satisfont
\[
 |X-F_+|^2=(a_S-f_S\cos\theta)^2,\qquad
 |X-F_-|^2=(a_S+f_S\cos\theta)^2.
\]
Comme \(a_S>f_S\), leurs racines positives ont pour somme \(2a_S\).
Pour \(\eta<0\), échanger les coordonnées donne la même preuve.
Dans le cas circulaire, \(f_S=0\). Enfin, substituer
\((Q,P)=r(\cos\phi,\sin\phi)\) dans
\(Q^2/a_S^2+P^2/b_S^2=1\) donne
\eqref{ii:eq:ii-elipse-radio-polar}. Le coefficient de \(r^2\) est
compris entre \(a^{-2}\) et \(b^{-2}\), ce qui démontre les bornes et la
caractérisation de l'intérieur.
\end{proof}

\subsection{Phase paramétrique, angle polaire et aire d'action}
\label{ii:sec:ii-elipse-fase-area}

\begin{proposition}[Loi aréale et deux coordonnées angulaires]
\label{ii:prop:ii-elipse-area}
La phase \(\theta\) de \(\gamma\) et son angle polaire continu
\(\phi\), relevés à partir de zéro, satisfont
\begin{equation}
 \phi(\theta)=\arg(a_S\cos\theta+i b_S\sin\theta),\qquad
 \frac{d\phi}{d\theta}
 =\frac{a_Sb_S}{a_S^2\cos^2\theta+b_S^2\sin^2\theta}>0.
 \label{ii:eq:ii-elipse-fases}
\end{equation}
Dans une carte où les deux tangentes sont définies,
\(\tan\phi=(b_S/a_S)\tan\theta\). L'aire orientée balayée est
\begin{equation}
 \mathcal A(\theta)=\frac12\int_0^\theta(QP'-PQ')\,dt
 =\hbar\theta,\qquad
 \mathcal A(1)=\hbar,\qquad \mathcal A(2\pi)=h.
 \label{ii:eq:ii-elipse-area-fase}
\end{equation}
\end{proposition}

\begin{proof}
La dérivée de l'argument d'un vecteur non nul est
\((QP'-PQ')/(Q^2+P^2)\). Ici, le numérateur est
\(a_Sb_S=2\hbar\), ce qui démontre la dérivée et la monotonie.
Le quotient \(P/Q\) donne la relation des tangentes ; le relèvement
continu fixe les quadrants et le nombre de tours. L'intégrale aréale
vaut \(a_Sb_S\theta/2=\hbar\theta\), et
\(h=2\pi\hbar\) donne la dernière identité.
\end{proof}

La phase paramétrique ne coïncide pas en général avec l’angle polaire. Un radian
de phase balaie \(\hbar\), mais ne s’identifie pas pour autant à un arc
circulaire ni à un secteur d’angle polaire unitaire. En coordonnées polaires,
la même loi s’écrit
\(d\mathcal A=\rho_{\partial}(\phi)^2d\phi/2\).
Avec l’orientation de \(\gamma\) utilisée ici,
\(\oint P\,dQ=-h\) : l’aire positive est \(h\), tandis que
l’intégrale de cette forme de degré un conserve son signe. Cela concorde avec la
distinction entre les transports symplectique et antisymplectique de la
section précédente.

\begin{figure}[htbp]
 \centering
 
% BEGIN OWNER /Users/ruben/Documents/ChatGPT/jueces y controles/REVISION_ARGUMENTAL_APERTURAS_CIERRES_20260915/II/manuscript_es/figures/elipse_accion_II.tex
% Geometría calculada, no coordenadas focales elegidas a ojo.
% Caso algebraico de ilustración xi=1/3; no evaluación de las constantes HMT.
\begingroup
\pgfmathsetmacro{\elXi}{1/3}
\pgfmathsetmacro{\elEta}{0.5*ln((1+\elXi)/(1-\elXi))}
\pgfmathsetmacro{\elA}{exp(\elEta/2)}
\pgfmathsetmacro{\elB}{exp(-\elEta/2)}
\pgfmathsetmacro{\elF}{sqrt(\elA*\elA-\elB*\elB)}
\pgfmathsetmacro{\elTheta}{30}
\pgfmathsetmacro{\elX}{\elA*cos(\elTheta)}
\pgfmathsetmacro{\elY}{\elB*sin(\elTheta)}
\pgfmathsetmacro{\elPolar}{atan(\elY/\elX)}
\begin{tikzpicture}[x=2.18cm,y=2.18cm,font=\small,
 every node/.style={inner sep=2pt},>=Stealth]
 \begin{scope}
  \node[font=\bfseries] at (0,1.38) {Forme et orientation};
  \fill[accent!12] (0,0)--(\elA,0)
    plot[domain=0:30,samples=50] ({\elA*cos(\x)},{\elB*sin(\x)})--cycle;
  \draw[ink,thick] (0,0) ellipse[x radius=\elA,y radius=\elB];
  \draw[->,black!55] (-1.42,0)--(1.48,0) node[right]{\(u\)};
  \draw[->,black!55] (0,-1.05)--(0,1.08) node[above]{\(v\)};
  \draw[accent,thick,-{Stealth}] (\elA,0)
    plot[domain=0:30,samples=50] ({\elA*cos(\x)},{\elB*sin(\x)});
  \draw[accent,thick] (0,0)--(\elX,\elY);
  \node[below left,font=\scriptsize] at (0,0) {\(O\)};
  \draw[accent] (.35,0) arc[start angle=0,end angle=\elPolar,radius=.35];
  \node[accent] at (.47,.065) {\(\phi\)};
  \fill[accent] (\elX,\elY) circle[radius=1.4pt];
  \node[above right] at (\elX,\elY) {\(X(\theta)\)};
  \node[accent,align=center,font=\scriptsize] at (1.26,.18) {\(\theta=\pi/6\)};
  \draw[dashed,black!45] (-\elF,0)--(\elX,\elY);
  \draw[dashed,black!45] (\elF,0)--(\elX,\elY);
  \fill[ink] (-\elF,0) circle[radius=1.5pt] node[below]{\(F_-\)};
  \fill[ink] (\elF,0) circle[radius=1.5pt] node[below]{\(F_+\)};
  \draw[<->,black!60] (-\elA,-1.04)--(0,-1.04)
    node[midway,below,font=\scriptsize]{\(a_S/\sqrt{2\hbar}\)};
  \draw[<->,black!60] (-1.32,0)--(-1.32,\elB)
    node[midway,left,font=\scriptsize]{\(\frac{b_S}{\sqrt{2\hbar}}\)};
  \node[align=center,font=\scriptsize] at (0,-1.55)
    {\(\rho_{\partial}(\phi)/\sqrt{2\hbar}=|OX|\)\\
     \(f_S^2=a_S^2-b_S^2\), \quad aire totale \(h\)};
 \end{scope}
 \begin{scope}[shift={(3.65,0)}]
  \node[font=\bfseries] at (0,1.38) {Forme réciproque};
  \draw[densely dashed,black!25] (0,0) ellipse[x radius=\elA,y radius=\elB];
  \draw[ink,thick] (0,0) ellipse[x radius=\elB,y radius=\elA];
  \draw[->,black!55] (-1.28,0)--(1.38,0) node[right]{\(u\)};
  \draw[->,black!55] (0,-1.30)--(0,1.26) node[right]{\(v\)};
  \fill[ink] (0,\elF) circle[radius=1.5pt] node[right]{\(F'_+\)};
  \fill[ink] (0,-\elF) circle[radius=1.5pt] node[right]{\(F'_-\)};
  \node[align=center,font=\scriptsize] at (0,-1.55)
    {\(\eta\mapsto-\eta\), \quad \(a_S\leftrightarrow b_S\)\\
     \(R_{\rm el}\mapsto R_{\rm el}^{-1}\), \quad \(a_Sb_S=2\hbar\)};
 \end{scope}
 \draw[->,accent,thick] (1.60,.80)--(2.17,.80)
   node[midway,above,font=\scriptsize]{échange};
\end{tikzpicture}
\endgroup

% END OWNER

 \caption{Demi-axes, foyers et réciprocité de l'ellipse d'action dans les
 coordonnées normalisées \(u=Q/\sqrt{2\hbar}\),
 \(v=P/\sqrt{2\hbar}\). Le secteur marqué possède la phase paramétrique
 \(\theta=\pi/6\), distincte de son angle polaire \(\phi\) ; son aire
 d'action est \(\hbar\pi/6\). La figure utilise le cas algébrique
 \(\xi=1/3\) pour rendre les relations lisibles : elle n'attribue pas cette valeur
 à la sortie angulaire HMT. Les foyers sont calculés à partir des demi-axes.
 À droite, l'échange des axes conserve l'aire totale \(h\)
 et inverse le rayon modulaire adimensionnel.}
 \label{ii:fig:ii-elipse-calculada}
\end{figure}

Pour le cas de contrôle de la figure, \(\xi=1/3\), les coordonnées
normalisées sont exactement
\[
 \eta=\tfrac12\log2,\quad
 \frac{a_S}{\sqrt{2\hbar}}=2^{1/4}=1.1892071150\ldots,\quad
 \frac{b_S}{\sqrt{2\hbar}}=
 \frac{f_S}{\sqrt{2\hbar}}=R_{\rm el}=2^{-1/4}=0.8408964153\ldots.
\]
Ce sont des nombres de contrôle de non-régression du dessin, non des valeurs attribuées au générateur
angulaire. Le certificat local vérifie également la forme circulaire et les
orientations \(\xi=\pm1/3,\pm3/5\).

\subsection{Réciprocité de forme et distance intérieure}

La publication rectangulaire restitue
\begin{equation}
 R_{\rm el}=\sqrt{b_S/a_S}=e^{-\eta/2},\qquad
 R_{\rm el}^4=\frac{1-\xi}{1+\xi},\qquad
 \xi=\frac{1-R_{\rm el}^4}{1+R_{\rm el}^4}.
 \label{ii:eq:ii-elipse-reciproca}
\end{equation}
La preuve consiste à substituer
\(2\eta=\log((1+\xi)/(1-\xi))\). L'échange des axes
envoie \(\eta\) sur \(-\eta\), \(R_{\rm el}\) sur
\(R_{\rm el}^{-1}\) et \(\tau=iR_{\rm el}^2\) sur \(-1/\tau\).
Le produit \(a_Sb_S\) est conservé. Le rayon \(R_{\rm el}\) est adimensionnel ;
le rayon polaire \(\rho_{\partial}\) a pour unité la racine carrée d'action.

L'intérieur admet également la distance projective de Hilbert. Si
\(z_-,x,y,z_+\) sont ordonnés sur une corde de l'ellipse, on définit
\[
 d_H(x,y)=\frac12\log
 \frac{|y-z_-|\,|x-z_+|}{|x-z_-|\,|y-z_+|}.
\]
C'est une distance adimensionnelle : les unités se simplifient dans le birapport.
La normalisation \((Q,P)\mapsto(Q/a_S,P/b_S)\) envoie l'ellipse
sur le disque unité et préserve ce rapport, car les longueurs sur une
même droite sont multipliées par un facteur commun. Elle transporte ainsi la métrique
vers le modèle de Beltrami--Klein. Ni cette distance intérieure ni la distance
entre les foyers ne s'identifie à une longueur atomique.

\subsection{Rayon de Bohr : conséquence de l'échelle d'action}
\label{ii:sec:ii-elipse-bohr}

Dans la réalisation électrodynamique, l'échelle de Bohr s'exprime par
\(a_{\rm B}=\hbar/(m_ec\alpha)\) \cite{ii:codata2022}. Dans la chaîne
du présent article, cette lecture succède aux constructions de
l'action, du vide, de la structure fine et de la masse électronique. On reçoit ici
\(\hbar,m_e,c,\alpha>0\) dans une même réalisation dimensionnelle.
\begin{proposition}[Expression aréale de l'échelle de Bohr]
\label{ii:prop:ii-elipse-bohr}
Dans cette réalisation,
\begin{equation}
 a_{\rm B}=\frac{a_Sb_S}{2m_ec\alpha}
 =\frac{\mathcal A(2\pi)}{2\pi m_ec\alpha}
 =\frac{\bar\lambda_C(m_e)}{\alpha},\qquad
 \bar\lambda_C(m_e)=\frac{\hbar}{m_ec}.
 \label{ii:eq:ii-elipse-bohr}
\end{equation}
\end{proposition}
\begin{proof}
Les identités \(a_Sb_S=2\hbar\) et
\(\mathcal A(2\pi)=2\pi\hbar\), substituées dans l'expression
de Bohr, donnent les deux premiers membres. La définition de
\(\bar\lambda_C(m_e)\) donne le troisième.
\end{proof}

Le numérateur a la dimension d'une action et \(m_ec\) celle d'une quantité de mouvement ; le
quotient est une longueur. Cette composition utilise le produit des
demi-axes et les trois autres grandeurs : elle n'affirme pas que Bohr soit un demi-axe,
un foyer, une distance de Hilbert ou le rayon modulaire. Elle n'assimile pas non plus
la masse électronique au quantum circulaire d'ordre 108. La notation
\(a_{\rm B}\) évite de confondre cette longueur avec la normalisation
\(a_0\) de l'opérateur d'aire.

\subsection{Réalisation inductive--capacitive et continuité constitutive}
\label{ii:sec:ii-elipse-lc}

Le pont électrodynamique de l'ellipse conserve une carte supplémentaire :
une horloge \(\omega>0\) et une impédance de référence \(Z_*>0\).
La construction reprise du corpus définit
\begin{equation}
 L_\eta=\frac{Z_*}{\omega}e^{-\eta},\qquad
 C_\eta=\frac{1}{Z_*\omega}e^\eta,
 \qquad Z_\eta=\sqrt{L_\eta/C_\eta}.
 \label{ii:eq:ii-elipse-lc}
\end{equation}
Ici, \(C_\eta\) est la capacité électrique, distincte du contre-angle \(C^*\).
\begin{proposition}[Produit dynamique et quotient constitutif]
\label{ii:prop:ii-elipse-lc}
Dans la carte précédente,
\begin{equation}
 L_\eta C_\eta=\omega^{-2},\qquad
 \frac{Z_\eta}{Z_*}=e^{-\eta}
 =\frac{b_S}{a_S}=R_{\rm el}^{2}.
 \label{ii:eq:ii-elipse-impedancia}
\end{equation}
Avec \(Q=\sqrt{Z_*}\,q\), \(P=\Phi/\sqrt{Z_*}\), l'énergie
\(H_{LC}=q^2/(2C_\eta)+\Phi^2/(2L_\eta)\) prend la forme
\[
 H_{LC}=\frac\omega2(e^{-\eta}Q^2+e^\eta P^2).
\]
Sur \(\gamma(\theta)\), ses composantes électrique et magnétique sont
\(U_E=\hbar\omega\cos^2\theta\) et
\(U_B=\hbar\omega\sin^2\theta\) ; leur somme vaut \(\hbar\omega\).
\end{proposition}
\begin{proof}
Dans le produit de \eqref{ii:eq:ii-elipse-lc}, les exponentielles et
\(Z_*\) se simplifient ; son quotient vaut \(Z_*^2e^{-2\eta}\). Prendre la racine
positive et utiliser \eqref{ii:eq:ii-elipse-reciproca} donne
\eqref{ii:eq:ii-elipse-impedancia}. Substituer \(q=Q/\sqrt{Z_*}\)
et \(\Phi=\sqrt{Z_*}P\) dans \(H_{LC}\) donne sa forme équilibrée.
Enfin, les demi-axes de \(\gamma\) simplifient les facteurs
\(e^{\mp\eta}\), et \(\cos^2\theta+\sin^2\theta=1\).
\end{proof}

Ce pont distingue la forme, l'horloge et le transducteur. Le choix
du sens dynamique doit s'accorder avec la convention symplectique : pour
les équations \(\dot Q=\partial_PH\),
\(\dot P=-\partial_QH\), la paramétrisation précédente satisfait
\(\dot\theta=-\omega\). Les identités énergétiques ne dépendent pas
de l'inversion de ce parcours.

Le lecteur résolvant du vide agit sur les canaux déjà produits \(q_\pm=e^{-(A_{\rm rad}\pm C^*_{\rm rad})}\) et l’incidence \(90/120\). Ses réponses \(r_\pm=(1-q_\pm^{90})/(1-q_\pm^{120})\) expriment \(\widehat Z=r_-/r_+\) et \(\widehat c=(r_+r_-)^{-1}\). Ce sont des réponses spectrales, non des rayons géométriques. L’identité \eqref{ii:eq:ii-elipse-impedancia} fixe ici la réalisation LC ; son identification à l’impédance constitutive du vide conserve l’application entre les deux lecteurs et leurs cartes. Leurs valeurs ne sont pas égalées au seul motif qu’elles partagent l’antécédent angulaire.
